\section{Related Work}
\label{sec:related_work}
The use of inertial measurement units (IMUs) for gait analysis has gained traction in both research and clinical domains, because of their  portability, affordability, 
and ability to collect kinematic data in real-world environments. There has been a number of studies that have validated the feasibility and accuracy of IMU-based 
gait monitoring in healthy, aging and impaired populations. An example of this case is the research by Yoon \textit{et al.} \cite{ref20}, where a cross-sectional 
study was conducted using wearable IMUs to assess gait parameters in younger and older adults, confirming their reliability compared to traditional 
magnetometer-based systems. Similarly, systematic reviews, like the one by Prisco \textit{et al.} \cite{ref21} and Zhu \textit{et al.} \cite{ref22} have 
comprehensively examined the validity of wearable inertial sensors, highlighting the effectiveness of single- and multi-point IMU configurations for capturing 
clinically relevant gait metrics.

Many studies have focused on gait event detection—such as identifying heel strikes or toe-offs—using IMUs. 
Voisard \textit{et al.} \cite{ref23} developed an automatic gait event detection algorithm capable of operating accurately across healthy 
individuals and patients with moderate to severe impairments, including those recovering from stroke. Larsen \textit{et al.} \cite{ref24} extended this line 
of research by introducing a neural network-based approach optimized for smartphone-embedded IMUs, achieving high detection accuracy in real-world conditions.

Despite the technical progress, much of the literature focuses on algorithmic sophistication or system validation rather than clinical interpretability or 
rehabilitation deployment. Bhole and Gokhale \cite{ref25} proposed an IMU-based system for basic gait event analysis, while Gudmundsson and 
Runarsson \cite{ref26} evaluated IMU sensor integration with heuristic rules for movement phase detection. However, these studies often lack the subject-specific 
generalization or low-complexity pipelines required for real-time, embedded medical use.

Notably, few works have been designed specifically for stroke rehabilitation, where gait asymmetry and inter-subject variability are prevalent. 
Patterson \textit{et al.} \cite{ref27} quantified temporal gait asymmetries among community-ambulating stroke survivors, while Zhao \textit{et al.} 
\cite{ref28} presented a dual foot-mounted IMU system for rehabilitation assessment in patients with cerebral thrombosis. Their system used an 
inequality-constrained ZUPT-aided inertial navigation approach to estimate gait parameters. However, these approaches often rely on dense feature sets 
or require manual tuning, making them less suitable for scalable deployment.

Although these previous studies have laid the groundwork for IMU-based gait analysis, there remains a lack of lightweight, interpretable, and clinically 
deployable methods tailored specifically for stroke-affected populations. This work addresses that gap by proposing a minimal-feature, high-performance 
classification pipeline using bilateral gyroscope extrema—optimized for subject-level generalization and embedded rehabilitation contexts.
